\documentclass[12pt, letterpaper]{article}

% ── Paquetes ──────────────────────────────────────────────────────────────────
\usepackage[utf8]{inputenc}
\usepackage[T1]{fontenc}
\usepackage[spanish]{babel}
\usepackage{geometry}
\usepackage{graphicx}
\usepackage{hyperref}
\usepackage{listings}
\usepackage{xcolor}
\usepackage{booktabs}
\usepackage{tabularx}
\usepackage{array}
\usepackage{fancyhdr}
\usepackage{titlesec}
\usepackage{enumitem}
\usepackage{tikz}
\usepackage{float}
\usepackage{microtype}
\usepackage{parskip}
\usepackage{mdframed}
\usetikzlibrary{shapes.geometric, arrows.meta, positioning}

% ── Geometría de página ───────────────────────────────────────────────────────
\geometry{
    top=2.5cm,
    bottom=2.5cm,
    left=3cm,
    right=3cm
}

% ── Colores ───────────────────────────────────────────────────────────────────
\definecolor{cgsp-blue}{RGB}{15, 76, 129}
\definecolor{cgsp-gray}{RGB}{240, 242, 245}
\definecolor{cgsp-green}{RGB}{39, 174, 96}
\definecolor{cgsp-red}{RGB}{192, 57, 43}
\definecolor{cgsp-orange}{RGB}{211, 84, 0}
\definecolor{code-bg}{RGB}{248, 249, 250}
\definecolor{code-comment}{RGB}{106, 153, 85}
\definecolor{code-keyword}{RGB}{86, 156, 214}
\definecolor{code-string}{RGB}{206, 145, 120}

% ── Estilo de listings (bloques de código) ────────────────────────────────────
\lstdefinestyle{cgsp-proto}{
    backgroundcolor=\color{code-bg},
    basicstyle=\ttfamily\small,
    breaklines=true,
    breakatwhitespace=false,
    frame=single,
    framesep=4pt,
    rulecolor=\color{cgsp-blue!40},
    xleftmargin=8pt,
    xrightmargin=8pt,
    showstringspaces=false,
    tabsize=2,
    commentstyle=\color{code-comment},
    keywordstyle=\color{code-keyword}\bfseries,
    stringstyle=\color{code-string},
    numbers=none,
    captionpos=b,
    keepspaces=true
}

\lstset{style=cgsp-proto}

% ── Encabezado y pie de página ────────────────────────────────────────────────
\pagestyle{fancy}
\fancyhf{}
\fancyhead[L]{\small\textcolor{cgsp-blue}{\textbf{RFC-CGSP}}}
\fancyhead[C]{\small CyberGame Socket Protocol (CGSP)}
\fancyhead[R]{\small Abril 2026}
\fancyfoot[C]{\thepage}
\renewcommand{\headrulewidth}{0.5pt}
\renewcommand{\headrule}{\hbox to\headwidth{\color{cgsp-blue}\leaders\hrule height \headrulewidth\hfill}}

% ── Títulos de sección ────────────────────────────────────────────────────────
\titleformat{\section}
    {\large\bfseries\color{cgsp-blue}}
    {\thesection.}{0.5em}{}[\titlerule]

\titleformat{\subsection}
    {\normalsize\bfseries\color{cgsp-blue!80}}
    {\thesubsection.}{0.5em}{}

\titleformat{\subsubsection}
    {\normalsize\bfseries\color{cgsp-blue!60}}
    {\thesubsubsection.}{0.5em}{}

% ── Cuadro de nota ────────────────────────────────────────────────────────────
\newmdenv[
    backgroundcolor=cgsp-gray,
    linecolor=cgsp-blue,
    linewidth=2pt,
    leftline=true,
    rightline=false,
    topline=false,
    bottomline=false,
    innerleftmargin=10pt,
    innerrightmargin=10pt,
    innertopmargin=6pt,
    innerbottommargin=6pt
]{nota}

% ── Cuadro de advertencia (nuevo: para notas de seguridad críticas) ───────────
\newmdenv[
    backgroundcolor=cgsp-red!8,
    linecolor=cgsp-red,
    linewidth=2pt,
    leftline=true,
    rightline=false,
    topline=false,
    bottomline=false,
    innerleftmargin=10pt,
    innerrightmargin=10pt,
    innertopmargin=6pt,
    innerbottommargin=6pt
]{advertencia}

% ── Hipervínculos ─────────────────────────────────────────────────────────────
\hypersetup{
    colorlinks=true,
    linkcolor=cgsp-blue,
    urlcolor=cgsp-blue,
    citecolor=cgsp-blue
}

% ─────────────────────────────────────────────────────────────────────────────
%  DOCUMENTO
% ─────────────────────────────────────────────────────────────────────────────
\begin{document}

% ── Portada ───────────────────────────────────────────────────────────────────
\begin{titlepage}
    \centering
    \vspace*{1cm}

    {\small\texttt{Especificación de Protocolo}}\\[0.3cm]
    {\Huge\bfseries\color{cgsp-blue} RFC-CGSP}\\[0.5cm]
    {\large\bfseries CyberGame Socket Protocol (CGSP)}

    \vspace{1.5cm}
    \rule{\textwidth}{1pt}\\[0.4cm]
    {\normalsize
        \begin{tabular}{ll}
            \textbf{Fecha:}    & Abril 2026 \\[4pt]
            \textbf{Materia:}  & Internet: Arquitectura y Protocolos \\[4pt]
            \textbf{Autores:}  & Andres Velez, Sebastian Salazar, Nathalia Cardoza
        \end{tabular}
    }\\[0.4cm]
    \rule{\textwidth}{1pt}

    \vspace{1.5cm}
    \begin{mdframed}[backgroundcolor=cgsp-gray, linecolor=cgsp-blue, linewidth=1pt]
        \textbf{Abstract.} Este documento especifica el \textit{CyberGame Socket Protocol
        (CGSP)}, un protocolo de la capa de aplicación diseñado para soportar la
        comunicación en tiempo real entre clientes y el servidor central de un simulador
        interactivo multijugador de ciberseguridad. El protocolo define el formato de
        mensajes, el vocabulario de comandos, las reglas de procedimiento y el manejo de
        errores para una partida en la que jugadores con roles de ``Atacante'' o
        ``Defensor'' interactúan sobre un plano virtual compartido.
    \end{mdframed}

    \vfill
    {\small Universidad EAFIT --- Departamento de Ingeniería de Sistemas}
\end{titlepage}

% ── Tabla de contenidos ───────────────────────────────────────────────────────
\tableofcontents
\newpage

% ═════════════════════════════════════════════════════════════════════════════
\section{Visión General del Protocolo}
% ═════════════════════════════════════════════════════════════════════════════

\subsection{Propósito}

CGSP permite que múltiples jugadores interactúen simultáneamente con un servidor de
juego central. A través de este protocolo, los jugadores pueden:

\begin{itemize}[leftmargin=2em]
    \item Autenticarse en el sistema mediante un servicio de identidad externo.
    \item Listar y unirse a salas de juego activas.
    \item Moverse por un plano virtual bidimensional.
    \item Ejecutar acciones específicas de su rol (atacar o defender recursos críticos).
    \item Recibir notificaciones de eventos del sistema en tiempo real.
\end{itemize}

\subsection{Modelo de Funcionamiento}

CGSP adopta el modelo \textbf{cliente-servidor} con comunicación \textbf{bidireccional
persistente}. El siguiente diagrama ilustra la arquitectura del sistema:

\begin{figure}[H]
\centering
\resizebox{\textwidth}{!}{%
\begin{tikzpicture}[
    node distance=2.2cm,
    box/.style={rectangle, draw=cgsp-blue, fill=cgsp-gray, rounded corners=4pt,
                minimum width=2.6cm, minimum height=1cm, align=center, font=\small},
    arrow/.style={-{Stealth[length=6pt]}, thick, color=cgsp-blue}
]
    \node[box] (webclient) {Navegador\\Web};
    \node[box, right=of webclient] (httpserver) {Servidor HTTP\\Python (8080)};
    \node[box, right=of httpserver] (servidor) {Servidor\\Central C (8081)};
    \node[box, right=of servidor] (identity) {Servicio de\\Identidad (9090)};
    \node[box, below=1.8cm of servidor] (clientes) {Clientes Desktop\\(Python/Java)};

    \draw[arrow] ([yshift=3pt]webclient.east) -- ([yshift=3pt]httpserver.west)
        node[midway, above, font=\tiny] {HTTP};
    \draw[arrow] ([yshift=-3pt]httpserver.west) -- ([yshift=-3pt]webclient.east)
        node[midway, below, font=\tiny] {HTTP};

    \draw[arrow] ([yshift=3pt]httpserver.east) -- ([yshift=3pt]servidor.west)
        node[midway, above, font=\tiny] {proxy CGSP};
    \draw[arrow] ([yshift=-3pt]servidor.west) -- ([yshift=-3pt]httpserver.east)
        node[midway, below, font=\tiny] {CGSP resp};

    \draw[arrow] ([yshift=3pt]servidor.east) -- ([yshift=3pt]identity.west)
        node[midway, above, font=\tiny] {TCP req/resp};
    \draw[arrow] ([yshift=-3pt]identity.west) -- ([yshift=-3pt]servidor.east)
        node[midway, below, font=\tiny] {OK ROLE / ERR};

    \draw[arrow] (clientes.north) -- (servidor.south)
        node[midway, right, font=\tiny] {TCP CGSP};
\end{tikzpicture}
}%
\caption{Arquitectura del sistema CGSP}
\end{figure}

El servidor acepta conexiones TCP entrantes y crea un hilo POSIX (\texttt{pthread})
dedicado por cada cliente. La conexión se mantiene activa durante toda la sesión del
jugador. El servidor también actúa como cliente frente al Servicio de Identidad para
autenticar usuarios. Un \textbf{Servidor HTTP Python} actúa como proxy entre el
navegador web y el servidor CGSP; los clientes desktop (Python/Tkinter para atacantes,
Java/Swing para defensores) se conectan directamente al servidor C.

\subsection{Capa de la Arquitectura TCP/IP}

CGSP opera en la \textbf{capa de aplicación} del modelo TCP/IP:

\begin{table}[H]
\centering
\begin{tabular}{>{\bfseries}l l}
\toprule
Capa & Protocolo \\
\midrule
Aplicación  & \textbf{CGSP}\\
Transporte  & TCP (SOCK\_STREAM) \\
Red         & IPv4 \\
Enlace      & Ethernet / Wi-Fi \\
\bottomrule
\end{tabular}
\caption{Pila de protocolos de CGSP}
\end{table}

\begin{nota}
\textbf{Justificación de TCP sobre UDP.} El juego requiere confiabilidad absoluta en
la entrega de mensajes. Un paquete de \texttt{ATTACK} o \texttt{DEFEND} que se
pierda comprometería la consistencia del estado del juego entre servidor y clientes.
TCP garantiza entrega ordenada, detección de pérdidas y retransmisión automática.
UDP sería adecuado solo para actualizaciones de posición de alta frecuencia donde
una pérdida es aceptable --- lo cual no aplica a este diseño.
\end{nota}

\subsection{Resolución de Nombres}

CGSP prohíbe el uso de direcciones IP codificadas (\textit{hardcoded}) en el
software. Toda referencia a servicios externos (servidor de juego, servicio de
identidad) se resuelve mediante DNS usando la función \texttt{getaddrinfo()} de
POSIX. Si la resolución falla, el servicio maneja la excepción sin finalizar su
ejecución, respondiendo al cliente con \texttt{ERR 503}.

% ═════════════════════════════════════════════════════════════════════════════
\section{Especificación del Servicio}
% ═════════════════════════════════════════════════════════════════════════════

\subsection{Primitivas del Servicio}

Las primitivas de CGSP se organizan según el estado del cliente en el servidor.

\subsubsection{Estado: CONECTADO (antes de autenticarse)}

\begin{table}[H]
\centering
\begin{tabularx}{\textwidth}{l l X}
\toprule
\textbf{Primitiva} & \textbf{Dirección} & \textbf{Descripción} \\
\midrule
\texttt{AUTH} & Cliente $\rightarrow$ Servidor & Autenticar un usuario contra el servicio de identidad \\
\texttt{QUIT} & Cliente $\rightarrow$ Servidor & Cerrar la conexión ordenadamente \\
\bottomrule
\end{tabularx}
\caption{Primitivas disponibles en estado CONECTADO}
\end{table}

\subsubsection{Estado: AUTENTICADO}

\begin{table}[H]
\centering
\begin{tabularx}{\textwidth}{l l X}
\toprule
\textbf{Primitiva} & \textbf{Dirección} & \textbf{Descripción} \\
\midrule
\texttt{LIST\_ROOMS}  & Cliente $\rightarrow$ Servidor & Obtener lista de salas activas \\
\texttt{CREATE\_ROOM} & Cliente $\rightarrow$ Servidor & Crear una nueva sala de juego \\
\texttt{JOIN}         & Cliente $\rightarrow$ Servidor & Unirse a una sala existente \\
\texttt{STATUS}       & Cliente $\rightarrow$ Servidor & Consultar estado actual del jugador \\
\texttt{QUIT}         & Cliente $\rightarrow$ Servidor & Cerrar la conexión \\
\bottomrule
\end{tabularx}
\caption{Primitivas disponibles en estado AUTENTICADO}
\end{table}

\subsubsection{Estado: EN SALA (WAITING)}

El jugador ha ejecutado \texttt{JOIN} exitosamente y se encuentra dentro de una sala
en espera de que la partida inicie. En este estado el servidor aún no ha recibido un
\texttt{START} válido (es decir, con al menos un atacante y un defensor presentes).

\begin{table}[H]
\centering
\begin{tabularx}{\textwidth}{l l X}
\toprule
\textbf{Primitiva} & \textbf{Dirección} & \textbf{Descripción} \\
\midrule
\texttt{START}  & Cliente $\rightarrow$ Servidor & Solicitar el inicio de la partida (requiere $\geq$1 atacante y $\geq$1 defensor) \\
\texttt{STATUS} & Cliente $\rightarrow$ Servidor & Consultar estado actual del jugador y la sala \\
\texttt{QUIT}   & Cliente $\rightarrow$ Servidor & Abandonar la sala y cerrar la conexión \\
\bottomrule
\end{tabularx}
\caption{Primitivas disponibles en estado EN SALA}
\end{table}

\subsubsection{Estado: EN PARTIDA (RUNNING)}

\begin{table}[H]
\centering
\begin{tabularx}{\textwidth}{l l l X}
\toprule
\textbf{Primitiva} & \textbf{Dirección} & \textbf{Rol} & \textbf{Descripción} \\
\midrule
\texttt{MOVE}   & C $\rightarrow$ S & Ambos    & Mover al jugador en el plano \\
\texttt{SCAN}   & C $\rightarrow$ S & ATTACKER & Detectar recursos críticos cercanos (radio 12) \\
\texttt{ATTACK} & C $\rightarrow$ S & ATTACKER & Lanzar un ataque sobre un recurso (debe estar a $\leq$12 unidades) \\
\texttt{DEFEND} & C $\rightarrow$ S & DEFENDER & Mitigar un ataque activo (debe estar a $\leq$12 unidades) \\
\texttt{STATUS} & C $\rightarrow$ S & Ambos    & Consultar estado del jugador/sala \\
\texttt{QUIT}   & C $\rightarrow$ S & Ambos    & Abandonar la partida \\
\bottomrule
\end{tabularx}
\caption{Primitivas disponibles en estado EN PARTIDA}
\end{table}

\subsubsection{Mensajes Asíncronos del Servidor}

\begin{table}[H]
\centering
\begin{tabularx}{\textwidth}{l l X}
\toprule
\textbf{Evento} & \textbf{Receptores} & \textbf{Descripción} \\
\midrule
\texttt{EVENT PLAYER\_JOINED}  & Sala completa    & Un nuevo jugador se unió a la sala \\
\texttt{EVENT PLAYER\_LEFT}    & Sala completa    & Un jugador abandonó la sala \\
\texttt{EVENT GAME\_STARTED}   & Sala completa    & La partida comenzó \\
\texttt{EVENT RESOURCE\_INFO}  & Solo defensores  & Coordenadas de recursos críticos (al iniciar) \\
\texttt{EVENT RESOURCE\_FOUND} & Atacantes de la sala & Recurso detectado por SCAN \\
\texttt{EVENT ATTACK}          & Sala completa    & Notificación de ataque activo \\
\texttt{EVENT DEFENDED}        & Sala completa    & Recurso fue defendido exitosamente \\
\texttt{EVENT ATTACK\_TIMEOUT} & Sala completa    & Recurso comprometido por expiración del timer \\
\texttt{EVENT GAME\_OVER}      & Sala completa    & Fin de partida con resultado \\
\bottomrule
\end{tabularx}
\caption{Mensajes asíncronos del servidor (broadcast y unicast)}
\end{table}

% ═════════════════════════════════════════════════════════════════════════════
\section{Formato de Mensajes}
% ═════════════════════════════════════════════════════════════════════════════

\subsection{Codificación General}

Todo mensaje CGSP debe cumplir las siguientes reglas de codificación:

\begin{itemize}[leftmargin=2em]
    \item \textbf{Codificación de caracteres:} UTF-8.
    \item \textbf{Delimitador de mensaje:} \texttt{\textbackslash{}n} (LF, \texttt{0x0A}). Un mensaje por línea.
    \item \textbf{Separador de campos:} espacio (\texttt{0x20}).
    \item \textbf{Longitud máxima de mensaje:} 2\,048 bytes.
    \item \textbf{Sensibilidad a mayúsculas en comandos:} NO (el servidor normaliza a mayúsculas).
    \item \textbf{Timeout de inactividad:} 600 segundos por defecto (configurable por \texttt{CGSP\_IDLE\_TIMEOUT}). El servidor enviará \texttt{ERR 408 Timeout: conexion cerrada por inactividad} y cerrará la conexión si el cliente no envía ningún mensaje en ese intervalo.
\end{itemize}

\textbf{Estructura general:}
\begin{lstlisting}
COMANDO [PARAM1] [PARAM2] ... [PARAMn]\n
\end{lstlisting}

\subsection{Mensajes del Cliente hacia el Servidor}

\subsubsection{AUTH --- Autenticación}

\begin{lstlisting}
AUTH <usuario> <password>\n
\end{lstlisting}

\begin{table}[H]
\centering
\begin{tabularx}{\textwidth}{l l l X}
\toprule
\textbf{Campo} & \textbf{Tipo} & \textbf{Máx.} & \textbf{Descripción} \\
\midrule
\texttt{usuario}  & string & 63 chars & Nombre de usuario registrado \\
\texttt{password} & string & 63 chars & Contraseña del usuario \\
\bottomrule
\end{tabularx}
\end{table}

\textbf{Ejemplo:} \texttt{AUTH carlos hack2026}

\subsubsection{LIST\_ROOMS --- Listar Salas}

\begin{lstlisting}
LIST_ROOMS\n
\end{lstlisting}

Sin parámetros. Solicita la lista de salas de juego activas.

\begin{nota}
\textbf{Información extendida de sala.} La respuesta \texttt{ROOM\_LIST} incluye ahora
el desglose de roles por sala (atacantes/defensores actuales), lo que permite al cliente
tomar decisiones de balanceo antes de unirse. Ver sección~\ref{sec:room-list}.
\end{nota}

\subsubsection{CREATE\_ROOM --- Crear Sala}

\begin{lstlisting}
CREATE_ROOM\n
\end{lstlisting}

Sin parámetros. Crea una nueva sala de juego y coloca los recursos críticos en posiciones aleatorias del mapa.

\subsubsection{JOIN --- Unirse a Sala}

\begin{lstlisting}
JOIN <room_id>\n
\end{lstlisting}

\begin{table}[H]
\centering
\begin{tabularx}{\textwidth}{l l X}
\toprule
\textbf{Campo} & \textbf{Tipo} & \textbf{Descripción} \\
\midrule
\texttt{room\_id} & entero positivo & Identificador numérico de la sala \\
\bottomrule
\end{tabularx}
\end{table}

\textbf{Ejemplo:} \texttt{JOIN 42}

\subsubsection{START --- Iniciar Partida}

\begin{lstlisting}
START\n
\end{lstlisting}

Solicita el inicio de la partida. Solo tiene efecto si hay al menos un atacante y un defensor en la sala.

\subsubsection{MOVE --- Mover Jugador}

\begin{lstlisting}
MOVE <dx> <dy>\n
\end{lstlisting}

\begin{table}[H]
\centering
\begin{tabularx}{\textwidth}{l l l X}
\toprule
\textbf{Campo} & \textbf{Tipo} & \textbf{Rango} & \textbf{Descripción} \\
\midrule
\texttt{dx} & entero & {[}$-100$, $100${]} & Desplazamiento en el eje X \\
\texttt{dy} & entero & {[}$-100$, $100${]} & Desplazamiento en el eje Y \\
\bottomrule
\end{tabularx}
\end{table}

El servidor clampea la posición resultante dentro del mapa $[0, 99] \times [0, 99]$. \\
\textbf{Ejemplo:} \texttt{MOVE 5 -3}

\subsubsection{SCAN --- Escanear Área}

\begin{lstlisting}
SCAN\n
\end{lstlisting}

Solo para jugadores con rol \texttt{ATTACKER}. Detecta recursos críticos en un radio de \textbf{12 unidades} (\texttt{DETECTION\_RADIUS}) alrededor de la posición actual del atacante. La detección se calcula con distancia euclidiana:
\[
    d = \sqrt{(x_\text{jugador} - x_\text{recurso})^2 + (y_\text{jugador} - y_\text{recurso})^2} \leq 12
\]
Si no se detecta ningún recurso, el servidor responde: \texttt{OK SCAN: no se encontraron recursos cercanos}.

\subsubsection{ATTACK --- Atacar Recurso}

\begin{lstlisting}
ATTACK <resource_id>\n
\end{lstlisting}

\begin{table}[H]
\centering
\begin{tabularx}{\textwidth}{l l X}
\toprule
\textbf{Campo} & \textbf{Tipo} & \textbf{Descripción} \\
\midrule
\texttt{resource\_id} & entero (0 ó 1) & Identificador del recurso crítico a atacar \\
\bottomrule
\end{tabularx}
\end{table}

\begin{nota}
\textbf{Restricción de proximidad.} El atacante debe estar a una distancia $\leq$ \texttt{DETECTION\_RADIUS} (12 unidades) del recurso para poder atacarlo. Si está demasiado lejos, el servidor responde \texttt{ERR 409 Debes estar cerca del recurso para atacarlo}.
\end{nota}

\textbf{Ejemplo:} \texttt{ATTACK 0}

\subsubsection{DEFEND --- Defender Recurso}

\begin{lstlisting}
DEFEND <resource_id>\n
\end{lstlisting}

\begin{table}[H]
\centering
\begin{tabularx}{\textwidth}{l l X}
\toprule
\textbf{Campo} & \textbf{Tipo} & \textbf{Descripción} \\
\midrule
\texttt{resource\_id} & entero (0 ó 1) & Identificador del recurso a defender \\
\bottomrule
\end{tabularx}
\end{table}

\begin{nota}
\textbf{Restricción de proximidad.} El defensor debe estar a una distancia $\leq$ \texttt{DETECTION\_RADIUS} (12 unidades) del recurso para poder defenderlo. Si está demasiado lejos, el servidor responde \texttt{ERR 409 Debes estar cerca del recurso para defenderlo}.
\end{nota}

\textbf{Ejemplo:} \texttt{DEFEND 0}

\subsubsection{STATUS --- Consultar Estado}

\begin{lstlisting}
STATUS\n
\end{lstlisting}

\subsubsection{QUIT --- Cerrar Sesión}

\begin{lstlisting}
QUIT\n
\end{lstlisting}

\subsection{Mensajes del Servidor hacia el Cliente}

\subsubsection{OK --- Operación Exitosa}

\begin{lstlisting}
OK <mensaje>\n
\end{lstlisting}

\textbf{Ejemplos:}
\begin{lstlisting}
OK Bienvenido carlos
OK Posicion: 15 42
OK Ataque lanzado
OK Recurso defendido
OK Hasta luego
\end{lstlisting}

\subsubsection{ROLE --- Rol Asignado}

\begin{lstlisting}
ROLE <ATTACKER|DEFENDER>\n
\end{lstlisting}

Enviado inmediatamente después de un \texttt{AUTH} exitoso.

\subsubsection{ERR --- Error}

\begin{lstlisting}
ERR <codigo> <descripcion>\n
\end{lstlisting}

\begin{table}[H]
\centering
\begin{tabularx}{\textwidth}{l l X}
\toprule
\textbf{Código} & \textbf{Nombre} & \textbf{Descripción} \\
\midrule
\texttt{400} & BAD\_REQUEST & Formato de mensaje inválido o parámetros faltantes \\
\texttt{401} & UNAUTHORIZED & Autenticación fallida o acción sin autenticar previamente \\
\texttt{403} & FORBIDDEN & El rol del jugador no permite esta acción \\
\texttt{404} & NOT\_FOUND & Sala o recurso no encontrado \\
\texttt{408} & TIMEOUT & Conexión cerrada por inactividad \\
\texttt{409} & CONFLICT & Conflicto de estado (sala llena, partida iniciada, recurso ya atacado) \\
\texttt{500} & INTERNAL\_ERROR & Error interno del servidor \\
\texttt{503} & SERVICE\_UNAVAILABLE & Servicio de identidad no disponible \\
\bottomrule
\end{tabularx}
\caption{Tabla de códigos de error de CGSP}
\end{table}

\textbf{Ejemplos:}
\begin{lstlisting}
ERR 401 Credenciales incorrectas
ERR 403 Solo los atacantes pueden usar ATTACK
ERR 404 Sala no encontrada
ERR 408 Timeout: conexion cerrada por inactividad
ERR 409 Recurso ya bajo ataque
ERR 503 Servicio de identidad no disponible, intenta luego
\end{lstlisting}

\subsubsection{ROOM\_CREATED --- Sala Creada}

\begin{lstlisting}
ROOM_CREATED <room_id>\n
\end{lstlisting}

\textbf{Ejemplo:} \texttt{ROOM\_CREATED 7}

\subsubsection{ROOM\_LIST --- Lista de Salas}
\label{sec:room-list}

\begin{lstlisting}
ROOM_LIST <n>\n
ROOM <id> <WAITING|RUNNING|FINISHED> <atacantes>/<defensores>/<max>\n
...
\end{lstlisting}

El campo de ocupación ahora desglosa \texttt{<atacantes>/<defensores>/<max>}, permitiendo al cliente evaluar el balance de roles antes de unirse.

\textbf{Ejemplo:}
\begin{lstlisting}
ROOM_LIST 2
ROOM 1 WAITING 1/0/10
ROOM 5 RUNNING 2/3/10
\end{lstlisting}

\subsubsection{Eventos del Servidor}

\begin{lstlisting}
EVENT PLAYER_JOINED <usuario> <ATTACKER|DEFENDER>\n
EVENT PLAYER_LEFT <usuario>\n
EVENT GAME_STARTED\n
EVENT RESOURCE_INFO <id> <x> <y>\n
EVENT RESOURCE_FOUND <id> <x> <y>\n
EVENT ATTACK <resource_id> <atacante>\n
EVENT DEFENDED <resource_id> <defensor>\n
EVENT ATTACK_TIMEOUT <resource_id>\n
EVENT GAME_OVER <ATTACKER|DEFENDER>\n
\end{lstlisting}

\begin{nota}
\textbf{Unicast vs.\ Broadcast.} Los eventos \texttt{RESOURCE\_INFO} se envían
únicamente a jugadores con rol defensor cuando inicia la partida. Los eventos \texttt{RESOURCE\_FOUND}
se envían al atacante que ejecutó el \texttt{SCAN} y también a los demás atacantes
de la misma sala. Todos los demás
eventos se difunden a todos los jugadores de la sala (\textit{broadcast}).
\end{nota}

% ═════════════════════════════════════════════════════════════════════════════
\section{Protocolo del Servicio de Identidad}
% ═════════════════════════════════════════════════════════════════════════════

\label{sec:identity}

El Servicio de Identidad es un componente externo al que el servidor CGSP consulta para
verificar credenciales de usuario. Esta sección especifica el sub-protocolo de
comunicación entre el servidor CGSP y dicho servicio.

\subsection{Modelo de Comunicación}

El servidor CGSP actúa como \textbf{cliente TCP} frente al Servicio de Identidad.
La conexión es de tipo \textbf{petición--respuesta}: se establece una conexión por
cada autenticación, se envía una consulta, se recibe la respuesta, y la conexión
se cierra.

\subsection{Formato de Mensajes}

Los mensajes siguen la misma codificación que CGSP: texto UTF-8 delimitado por
\texttt{\textbackslash{}n}.

\textbf{Petición del servidor CGSP al Servicio de Identidad:}
\begin{lstlisting}
AUTH_CHECK <usuario> <password>\n
\end{lstlisting}

\textbf{Respuestas del Servicio de Identidad:}
\begin{lstlisting}
OK <ATTACKER|DEFENDER>\n   % Credenciales validas; retorna el rol asignado
ERR 401 <descripcion>\n    % Credenciales invalidas
ERR 503 <descripcion>\n    % Servicio temporalmente no disponible
\end{lstlisting}

\begin{nota}
El rol del jugador (\texttt{ATTACKER} o \texttt{DEFENDER}) es devuelto directamente
por el Servicio de Identidad. El servidor CGSP no asigna roles de forma independiente,
garantizando así consistencia entre partidas.
\end{nota}

\subsection{Ejemplo de Intercambio}

\begin{lstlisting}
# Servidor CGSP -> Servicio de Identidad
AUTH_CHECK carlos hack2026

# Servicio de Identidad -> Servidor CGSP  (credenciales OK)
OK ATTACKER

# Servicio de Identidad -> Servidor CGSP  (credenciales incorrectas)
ERR 401 Usuario o contrasena invalidos

# Servicio de Identidad -> Servidor CGSP  (servicio caido)
ERR 503 Base de datos no disponible
\end{lstlisting}

\subsection{Configuración}

El host y puerto del Servicio de Identidad se configuran mediante variables de
entorno (ver sección~6). El servidor de juego nunca almacena credenciales ni resultados
en caché; cada \texttt{AUTH} del cliente genera una nueva consulta al Servicio de
Identidad.

% ═════════════════════════════════════════════════════════════════════════════
\section{Reglas de Procedimiento (Flujo de Operación)}
% ═════════════════════════════════════════════════════════════════════════════

\subsection{Máquina de Estados del Cliente}

\begin{figure}[H]
\centering
\begin{tikzpicture}[
    node distance=2.2cm,
    state/.style={rectangle, draw=cgsp-blue, fill=cgsp-gray, rounded corners=6pt,
                  minimum width=5.5cm, minimum height=1cm, align=center, font=\small\bfseries},
    arrow/.style={-{Stealth[length=6pt]}, thick, color=cgsp-blue!70},
    label/.style={font=\tiny, color=cgsp-blue!80, align=center}
]
    \node[state] (S1) {CONECTADO\\{\footnotesize AUTH, QUIT}};
    \node[state, below=of S1] (S2) {AUTENTICADO\\{\footnotesize LIST\_ROOMS, CREATE\_ROOM, JOIN, STATUS, QUIT}};
    \node[state, below=of S2] (S3) {EN SALA (WAITING)\\{\footnotesize START, STATUS, QUIT}};
    \node[state, below=of S3] (S4) {EN PARTIDA (RUNNING)\\{\footnotesize MOVE, SCAN/ATTACK, DEFEND, STATUS, QUIT}};

    \draw[arrow] (S1) -- (S2)
        node[midway, right, label] {AUTH exitoso\\$\rightarrow$ OK + ROLE};
    \draw[arrow] (S2) -- (S3)
        node[midway, right, label] {JOIN exitoso\\$\rightarrow$ OK + EVENT PLAYER\_JOINED (broadcast)};
    \draw[arrow] (S3) -- (S4)
        node[midway, right, label] {START ($\geq$1 ATC + $\geq$1 DEF)\\$\rightarrow$ EVENT GAME\_STARTED + EVENT RESOURCE\_INFO (defensores)};
\end{tikzpicture}
\caption{Máquina de estados del cliente CGSP}
\end{figure}

\begin{nota}
\textbf{Posiciones de spawn.} Al unirse a una sala, el atacante comienza en la posición $(0,\,0)$ y el defensor en $(99,\,99)$ (esquinas opuestas del mapa $100\times100$). Esto garantiza que el atacante deba explorar el mapa para encontrar recursos.
\end{nota}

\subsection{Ciclo de Vida de una Sala}

\begin{figure}[H]
\centering
\begin{tikzpicture}[
    node distance=3.5cm,
    state/.style={rectangle, draw=cgsp-blue, fill=cgsp-gray, rounded corners=4pt,
                  minimum width=2.8cm, minimum height=0.9cm, align=center, font=\small\bfseries},
    arrow/.style={-{Stealth[length=6pt]}, thick, color=cgsp-blue!70}
]
    \node[state] (W) {WAITING};
    \node[state, right=of W] (R) {RUNNING};
    \node[state, right=of R] (F) {FINISHED};

    \draw[arrow] (W) -- (R) node[midway, above, font=\tiny] {START exitoso};
    \draw[arrow] (R) -- (F) node[midway, above, font=\tiny] {condición\\de victoria};
\end{tikzpicture}
\caption{Ciclo de vida de una sala de juego}
\end{figure}

\subsection{Condiciones de Victoria}
\label{sec:victoria}

La partida finaliza cuando se cumple \textbf{cualquiera} de las siguientes condiciones,
momento en que el servidor emite \texttt{EVENT GAME\_OVER <ganador>} a todos los
jugadores de la sala:

\begin{table}[H]
\centering
\begin{tabularx}{\textwidth}{l X}
\toprule
\textbf{Ganador} & \textbf{Condición} \\
\midrule
\texttt{ATTACKER} & Al menos \textbf{un recurso crítico} no es defendido dentro del plazo de 30 s tras ser atacado. El servidor emite \texttt{EVENT ATTACK\_TIMEOUT <resource\_id>} seguido de \texttt{EVENT GAME\_OVER ATTACKER} a todos los jugadores de la sala. \\[6pt]
\texttt{DEFENDER} & La sala lleva en estado \texttt{RUNNING} más de \texttt{CGSP\_GAME\_TIMEOUT} segundos (por defecto 600~s) sin que ningún ataque expire. El servidor emite \texttt{EVENT GAME\_OVER DEFENDER} automáticamente. \\
\bottomrule
\end{tabularx}
\caption{Condiciones de victoria de CGSP}
\end{table}

\begin{nota}
Si un atacante abandona la partida (\texttt{QUIT}) mientras hay un ataque en curso,
dicho ataque se cancela automáticamente y se emite \texttt{EVENT PLAYER\_LEFT}.
Si todos los atacantes abandonan, el servidor declara victoria del equipo defensor.
\end{nota}

\subsection{Flujo de Ataque y Defensa}

\begin{lstlisting}
ATACANTE          SERVIDOR                  DEFENSOR(ES)
   |                  |                          |
   |-- ATTACK 0 ----->|                          |  (atacante cerca del recurso)
   |<- OK Ataque lanzado                         |
   |                  |-- EVENT ATTACK 0 carlos->|  (broadcast a toda la sala)
   |                  |     [timer: 30 s]         |
   |                  |<-- DEFEND 0 -------------|  (defensor cerca del recurso)
   |                  |-- EVENT DEFENDED 0 maria->|  (broadcast a toda la sala)
   |<- EVENT DEFENDED 0 maria                    |
   |                  |                          |
  (Escenario alternativo: el timer de 30 s expira sin que nadie defienda)
   |                  |-- EVENT ATTACK_TIMEOUT 0->|  (broadcast)
   |<- EVENT ATTACK_TIMEOUT 0                    |
   |                  |-- EVENT GAME_OVER ATTACKER>|  (broadcast)
   |<- EVENT GAME_OVER ATTACKER                  |
\end{lstlisting}

\begin{nota}
\textbf{Timer de ataque.} Un recurso bajo ataque que no sea defendido dentro de
\textbf{30 segundos} constituye victoria para el equipo atacante. El servidor emite
\texttt{EVENT ATTACK\_TIMEOUT <resource\_id>} seguido de \texttt{EVENT GAME\_OVER ATTACKER}
a todos los jugadores de la sala.
\end{nota}

\subsection{Flujo de Exploración del Atacante}

\begin{lstlisting}
ATACANTE              SERVIDOR
   |                     |
   |-- MOVE 30 35 ----->|   pos = (30, 35)
   |<- OK Posicion: 30 35
   |-- SCAN ------------>|   dist((30,35), recurso) <= 12 ?
   |<- EVENT RESOURCE_FOUND 0 32 37    (si detecta)
   |<- OK SCAN: no se encontraron recursos cercanos  (si no detecta)
   |-- ATTACK 0 -------->|   dist((30,35), recurso_0) <= 12 ?
   |<- OK Ataque lanzado |   (si esta cerca)
   |<- ERR 409 Debes estar cerca ...   (si muy lejos)
\end{lstlisting}

\subsection{Manejo de Errores y Excepciones}

\begin{table}[H]
\centering
\begin{tabularx}{\textwidth}{X X}
\toprule
\textbf{Situación} & \textbf{Respuesta del Servidor} \\
\midrule
Línea vacía o sin \texttt{\textbackslash{}n} & Ignorada silenciosamente \\
Comando desconocido & \texttt{ERR 400 Comando desconocido} \\
Parámetros faltantes & \texttt{ERR 400 Uso: CMD <params>} \\
Acción sin autenticar & \texttt{ERR 401 Debes autenticarte primero} \\
Rol incorrecto para acción & \texttt{ERR 403 Solo los [rol] pueden usar [CMD]} \\
Sala inexistente & \texttt{ERR 404 Sala no encontrada} \\
Recurso inexistente & \texttt{ERR 404 Recurso no encontrado} \\
Sala llena o partida iniciada & \texttt{ERR 409 <descripción del conflicto>} \\
Servicio de identidad caído & \texttt{ERR 503} --- servidor \textbf{no} finaliza \\
Desconexión abrupta del cliente & \texttt{read()} retorna $\leq 0$ $\rightarrow$ cierre limpio del hilo \\
Inactividad $>$ 600 s & \texttt{ERR 408 Timeout: conexion cerrada por inactividad} $\rightarrow$ cierre del hilo \\
AUTH con sesión ya activa & \texttt{ERR 403 Ya estas autenticado} \\
JOIN cuando ya está en sala & \texttt{ERR 403 Ya estas en una sala} \\
\bottomrule
\end{tabularx}
\caption{Tabla de manejo de excepciones del protocolo}
\end{table}

% ═════════════════════════════════════════════════════════════════════════════
\section{Ejemplos de Implementación}
% ═════════════════════════════════════════════════════════════════════════════

\subsection{Sesión Completa de un Atacante}

\begin{lstlisting}
# Conexion establecida
S->C: OK Bienvenido al servidor CyberGame (CGSP v2.0). Usa AUTH

# Autenticacion
C->S: AUTH carlos hack2026
S->C: OK Bienvenido carlos
S->C: ROLE ATTACKER

# Listar salas disponibles
C->S: LIST_ROOMS
S->C: ROOM_LIST 1
S->C: ROOM 3 WAITING 0/1/10   # 0 atacantes, 1 defensor, max 10

# Unirse a sala
C->S: JOIN 3
S->C: OK Te uniste a la sala

# Iniciar partida
C->S: START
S->C: OK Partida iniciada
S->C: EVENT GAME_STARTED

# Explorar el mapa
C->S: MOVE 30 35
S->C: OK Posicion: 30 35

C->S: SCAN
S->C: EVENT RESOURCE_FOUND 0 32 37

# Lanzar ataque
C->S: ATTACK 0
S->C: OK Ataque lanzado

# Notificacion de defensa exitosa
S->C: EVENT DEFENDED 0 maria

C->S: QUIT
S->C: OK Hasta luego
\end{lstlisting}

\subsection{Sesión Completa de un Defensor}

\begin{lstlisting}
C->S: AUTH maria seg2026
S->C: OK Bienvenido maria
S->C: ROLE DEFENDER

C->S: JOIN 3
S->C: OK Te uniste a la sala

C->S: START
S->C: EVENT GAME_STARTED
S->C: EVENT RESOURCE_INFO 0 32 37    # conoce los recursos
S->C: EVENT RESOURCE_INFO 1 75 82

C->S: MOVE 32 37
S->C: OK Posicion: 32 37

# Alerta de ataque
S->C: EVENT ATTACK 0 carlos

C->S: DEFEND 0
S->C: OK Recurso defendido

C->S: QUIT
S->C: OK Hasta luego
\end{lstlisting}

\subsection{Escenario: Ataque no Defendido (victoria atacante)}

\begin{lstlisting}
# Timer de 30 s expira sin respuesta del defensor
S->C (ambos): EVENT ATTACK_TIMEOUT 0
S->C (ambos): EVENT GAME_OVER ATTACKER
\end{lstlisting}

\subsection{Diagrama de Secuencia de una Partida Completa}

\begin{lstlisting}
Atacante        Servidor         Defensor
   |               |                |
   |-- AUTH ------>|<----- AUTH ----|
   |<- OK/ROLE ----|---- OK/ROLE -->|
   |               |                |
   |-- CREATE_ROOM->|               |
   |<- ROOM_CREATED|                |
   |               |                |
   |-- JOIN 5 ---->|<----- JOIN 5 --|
   |<- OK ----------|----- OK ----->|
   |               |                |
   |-- START ------>|<----- START --|
   |<- GAME_STARTED |-- GAME_STARTED>|
   |                |-- RESOURCE_INFO>|
   |               |                |
   |-- MOVE/SCAN -->|               |
   |<- FOUND -------|               |
   |               |                |
   |-- ATTACK 0 -->|-- EVENT ATTACK>|
   |<- OK ---------|<--- DEFEND 0 --|
   |<- DEFENDED ----|---- OK ------>|
   |               |                |
   |       (o bien, timer expira)   |
   |<- ATTACK_TIMEOUT 0 ----------->|
   |<- GAME_OVER ATTACKER ---------->|
\end{lstlisting}

\subsection{Prueba Manual con Telnet}

Para verificar el protocolo sin necesidad de un cliente completo:

\begin{lstlisting}[language=bash]
# Terminal 1: Iniciar el servicio de identidad
cd Protocol_ServiceID
python3 identity_server.py 9090 identity.log

# Terminal 2: Iniciar el servidor de juego
export IDENTITY_HOST=localhost
export IDENTITY_PORT=9090
export CGSP_IDLE_TIMEOUT=600   # opcional
export CGSP_GAME_TIMEOUT=600   # opcional
cd server && ./server 8081 server.log

# Terminal 3 (opcional): Servidor HTTP/Lobby
export GAME_HOST=localhost
export GAME_PORT=8081
python3 cliente_web_y_juego/servidor_http/server_http.py
# Abrir http://127.0.0.1:8080 en el navegador

# Terminal 4: Conectar directamente como atacante
telnet localhost 8081

# Terminal 5: Conectar directamente como defensor
telnet localhost 8081
\end{lstlisting}

% ═════════════════════════════════════════════════════════════════════════════
\section{Parámetros del Servidor}
% ═════════════════════════════════════════════════════════════════════════════

\subsection*{Servidor Principal (C)}
\begin{table}[H]
\centering
\small
\begin{tabularx}{\textwidth}{>{\ttfamily\raggedright\arraybackslash}p{4.2cm} >{\raggedright\arraybackslash}X >{\raggedright\arraybackslash}p{2cm} >{\raggedright\arraybackslash}p{3.8cm}}
\toprule
\normalfont\textbf{Parámetro} & \textbf{Fuente} & \textbf{Default} & \textbf{Descripción} \\
\midrule
\normalfont Puerto de escucha  & \texttt{./server <p>}       & ---     & Puerto TCP del servidor CGSP \\
\normalfont Archivo de logs    & \texttt{./server <p> <log>} & ---     & Archivo de logging \\
IDENTITY\_HOST  & Variable de entorno         & localhost & Hostname del servicio de identidad \\
IDENTITY\_PORT  & Variable de entorno         & 9090      & Puerto del servicio de identidad \\
CGSP\_IDLE\_TIMEOUT & Variable de entorno      & 600 s     & Timeout de inactividad del cliente \\
CGSP\_GAME\_TIMEOUT & Variable de entorno      & 600 s     & Duración máxima de partida \\
\bottomrule
\end{tabularx}
\caption{Parámetros del servidor CGSP (C)}
\end{table}

\subsection*{Servicio de Identidad (\texttt{identity\_server.py})}
\begin{table}[H]
\centering
\small
\begin{tabularx}{\textwidth}{>{\ttfamily\raggedright\arraybackslash}p{4.2cm} >{\raggedright\arraybackslash}X >{\raggedright\arraybackslash}p{2.2cm} >{\raggedright\arraybackslash}p{3.8cm}}
\toprule
\normalfont\textbf{Parámetro} & \textbf{Fuente} & \textbf{Default} & \textbf{Descripción} \\
\midrule
\normalfont Puerto    & \texttt{identity\_server.py <p>}      & 9090         & Puerto de escucha \\
\normalfont Archivo de log & \texttt{identity\_server.py <p> <log>}& identity.log & Log del servicio \\
IDENTITY\_IDLE\_TIMEOUT & Variable de entorno & 120 s       & Timeout de inactividad por cliente \\
IDENTITY\_USERS\_FILE  & Variable de entorno & users.json   & Ruta al archivo de usuarios \\
\bottomrule
\end{tabularx}
\caption{Parámetros del servicio de identidad}
\end{table}

\subsection*{Servidor HTTP / Lobby y Clientes Desktop}
\begin{table}[H]
\centering
\small
\begin{tabularx}{\textwidth}{>{\ttfamily\raggedright\arraybackslash}p{3.2cm} >{\raggedright\arraybackslash}p{2cm} >{\raggedright\arraybackslash}X}
\toprule
\normalfont\textbf{Variable} & \textbf{Default} & \textbf{Descripción} \\
\midrule
GAME\_HOST  & localhost   & Host del servidor CGSP (servidor HTTP) \\
GAME\_PORT  & 8081        & Puerto del servidor CGSP (servidor HTTP) \\
HTTP\_HOST  & 127.0.0.1   & Interfaz de escucha del servidor HTTP \\
HTTP\_PORT  & 8080        & Puerto del servidor HTTP \\
CGSP\_HOST  & localhost   & Host del servidor CGSP (clientes desktop) \\
CGSP\_PORT  & 8081        & Puerto del servidor CGSP (clientes desktop) \\
CGSP\_USER  & atacante1   & Usuario para autenticación automática \\
CGSP\_PASS  & pass123     & Contraseña para autenticación automática \\
CGSP\_ROOM  & 1           & Sala a la que se une el cliente al arrancar \\
\bottomrule
\end{tabularx}
\caption{Variables de entorno del servidor HTTP y clientes desktop}
\end{table}

% ═════════════════════════════════════════════════════════════════════════════
\section{Parámetros del Mapa de Juego}
% ═════════════════════════════════════════════════════════════════════════════

\begin{table}[H]
\centering
\begin{tabular}{l r}
\toprule
\textbf{Parámetro} & \textbf{Valor} \\
\midrule
Dimensiones del mapa             & $100 \times 100$ unidades \\
Recursos críticos por sala       & 2 (posición aleatoria al crear sala) \\
Radio de detección (\textsc{SCAN}) & 12 unidades (distancia euclidiana) \\
Tiempo límite de defensa         & 30 segundos \\
Tiempo máximo de partida         & 600 segundos (10 minutos); victoria del defensor al expirar \\
Máx.\ jugadores por sala         & 10 \\
Máx.\ salas simultáneas          & 20 \\
\bottomrule
\end{tabular}
\caption{Parámetros del mapa del juego}
\end{table}

% ═════════════════════════════════════════════════════════════════════════════
\section{Consideraciones de Seguridad}
% ═════════════════════════════════════════════════════════════════════════════

\begin{itemize}[leftmargin=2em]
    \item Las contraseñas \textbf{no se almacenan} en el servidor de juego; se reenvían al
          servicio de identidad en cada autenticación.
    \item Las contraseñas \textbf{no se registran} en los logs del servidor (aparecen como
          \texttt{[***]}).
    \item Cada cliente opera en un hilo independiente; un cliente malicioso no puede
          bloquear a los demás.
        \item El timeout de inactividad de 600 s mitiga conexiones ``zombie'' que
          consuman hilos del servidor indefinidamente.
\end{itemize}

\begin{advertencia}
\textbf{Advertencia de seguridad --- Ausencia de cifrado.}
Esta versión del protocolo transmite todos los mensajes, incluyendo credenciales de
autenticación, en \textbf{texto plano sin cifrado}. Esto representa un riesgo real
de interceptación de contraseñas en redes no confiables.
\end{advertencia}

% ─────────────────────────────────────────────────────────────────────────────
\vspace{2cm}
\hrule
\vspace{0.3cm}
{\small\centering\textit{Fin del documento RFC-CGSP --- CyberGame Socket Protocol}\par}

\end{document}